% ============================================================================
\chapter{결론}\label{ch:conclusion}
% ============================================================================

\section{연구 요약}

본 연구는 갑상선 결절 초음파 영상에서 경계 정확도를 높이면서 CPU에서 실시간으로 동작하는 검출·분할
모델을 목표로, YOLOv11-seg를 기반으로 데이터 혼합, 경계 정교화, 사후 양자화를 하나의 파이프라인에서
검증하였다. 세 가설의 검증 결과는 다음과 같다.

\begin{enumerate}
  \item \textbf{경량 백본에서 출발해도 손해가 없다 (H1 지지).} YOLOv11n-seg는 small 대비 파라미터 1/3.5,
    연산량 1/3.4임에도 held-out Mask \mAP{}가 0.6272로 small(0.6239)과 동등하였다. 결절이라는 단일
    클래스 과제에서 큰 모델의 추가 용량은 정확도로 이어지지 않았다.

  \item \textbf{경계 정교화를 추론 비용 없이 얻을 수 있다 (H2 지지).} 학습 전용 경계 분기와 경계 가중 손실로
    Dice가 0.8195에서 0.8300으로, HD95가 9.628\%에서 8.378\%로(13.0\% 감소), Boundary IoU가 0.4358에서
    0.4414로 개선되었고, 세 지표 모두 영상 단위 대응표본 검정에서 유의하였다. 추론 연산량은 9.73 GFLOPs로
    베이스라인과 같았다. 다만 손실 계수를 held-out 결과로 고른 탓에 이 개선 폭은 낙관적일 수 있으며,
    검증 분할로 다시 점검했을 때 경계 개선 자체는 모든 계수 구성에서 재현되었다.

  \item \textbf{CPU 배포가 실용 범위에 들어왔다 (H3 지지, 변환 도구에 따라 다름).} OpenVINO INT8 모델은 3.71 MB,
    CPU 지연 4.10 ms로 ONNX FP32 대비 15.3배 빨랐고 Dice 손실은 0.001에 그쳤다. 반면 ONNX Runtime INT8은
    검출 헤드를 부동소수점으로 남겨야 해 44.00 ms에 머물렀고 Dice도 0.012 떨어졌다.
\end{enumerate}

교차 데이터셋 조건에서는 경계 개선이 유지되지 않았다. 다만 3개 시드로 반복하자 악화 역시 통계적으로
확인되지 않았으므로, 본 연구는 교차 데이터셋 일반화에 대해 판단을 유보한다. 단일 실행 결과만 보았다면
"교차 도메인에서 크게 악화된다"는 잘못된 결론을 내렸을 것이며, 이를 시드 반복으로 바로잡은 과정 자체도
본 연구가 남기는 교훈이다.

\section{향후 연구}

\paragraph{오검출 실패 유형의 억제.} 정성 분석의 최악 사례(\figref{fig:qualitative})처럼 제안 모델이
주변 조직을 결절로 포함하는 경우가 있었다. 경계 가중이 잘못 잡힌 경계를 오히려 강화했을 가능성이 있으므로,
예측 경계의 불확실성을 함께 추정하거나 오검출에 별도 벌점을 주는 손실을 검토할 필요가 있다.

\paragraph{도메인 이동 대응.} \chapref{ch:dataset}의 분석은 두 데이터셋의 차이가 영상 형식(장비 화면
캡처와 확대 영상), 결절 크기 분포, 결절 수에서 모두 크다는 것을 보여 준다. DDTI의 장비 UI 영역을
가리거나 잘라내는 전처리, 스타일 정규화, 도메인 적대 학습\citep{ganin2015dann}을 결합하여 교차 데이터셋
성능을 다시 검증할 필요가 있다. 또한 더 많은 시드와 제3의 외부 데이터셋을 확보하면 일반화 여부를
판정할 수 있는 검정력을 얻을 수 있다.

\paragraph{양자화 인식 학습.} ONNX Runtime INT8에서 나타난 정확도 손실은 사후 양자화의 한계일 수 있다.
양자화 인식 학습으로 헤드까지 정수 연산으로 옮기면서 정확도를 회복할 수 있는지 확인할 가치가 있다.

\paragraph{실제 배포 장비에서의 측정.} 본 연구의 CPU 지연은 96코어 서버 CPU에서 측정한 값이다.
휴대형 초음파나 판독 PC에 쓰이는 저전력 x86·ARM CPU, Jetson 같은 임베디드 장비에서 지연과 전력을
실측해야 배포 가능성을 확정할 수 있다.

\paragraph{임상적 유용성 검증.} 경계 지표의 개선이 판독에 의미 있는 차이인지는 별도로 확인해야 한다.
예측 경계로부터 TI-RADS 경계 범주(매끄러움, 분엽상, 불규칙)를 자동 판정하고, 그 일치도를 전문의 판독과
비교하는 연구로 확장할 수 있다.
